\documentclass{article}
\usepackage[utf8]{inputenc}
\usepackage[T1]{fontenc}
\usepackage{geometry}
\usepackage{graphicx}
\usepackage{float}
\usepackage{hyperref}
\usepackage[english,indonesian]{babel}

\title{Optimasi Pemrosesan Template LaTeX Menggunakan Pendekatan Berbasis AI}
\author{Academic Support Agent}
\date{\today}

\begin{document}
\maketitle

\begin{abstract}

\end{abstract}

\chapter{PENDAHULUAN}


\chapter{TINJAUAN PUSTAKA}


% Original content (commented out):
% % ----------------------------------------------------------
% Template para Trabalho de Conclusão de Curso (TCC)
% ----------------------------------------------------------
%
% Este é um template LaTeX desenvolvido para auxiliar estudantes na elaboração de seus Trabalhos de Conclusão de Curso (TCCs). O template segue as normas da ABNT (Associação Brasileira de Normas Técnicas) e já inclui diversas seções pré-configuradas, facilitando a organização e formatação do trabalho.
%
% Como utilizar?
%
% Dados do TCC: 
% Preencha o arquivo DadosDoTCC com as informações do seu trabalho, como título, autor, orientador, etc.
%
% Elementos Pré-Textuais:
% Inclua os elementos pré-textuais necessários, como a capa, folha de rosto, agradecimentos, resumos, listas de figuras, tabelas, siglas e símbolos.
%
% Sumário:
% O sumário é gerado automaticamente e inclui links para as seções do documento.
%
% Elementos Textuais:
% Inclua os capítulos principais do seu TCC, como introdução, objetivos, justificativa, revisão bibliográfica, metodologia, resultados, etc.
%
% Elementos Pós-Textuais:
% Inclua as referências bibliográficas, apêndices e anexos conforme necessário.
%
% Customização:
% Este template pode ser facilmente personalizado para atender às necessidades específicas do seu TCC. Você pode modificar o conteúdo dos arquivos de texto .tex e ajustar as configurações conforme desejar.
%
% Espero que este template facilite o processo de elaboração do seu TCC, permitindo que você se concentre no conteúdo do trabalho sem se preocupar com detalhes de formatação.
% ----------------------------------------------------------
% Elaborado por Julio Cesar Sampaio Dutra e  Wancley Oinhos Pedruzzi
% E-mail: julio.dutra@ufes.br; wancleypedruzzi@hotmail.com 
% Data: 05 de abril de 2024
% ----------------------------------------------------------
% Importar os pacots e configurações básicas para atender à ABNT e o template do curso
\input{NaoAlterar/configuracoes}
\makeindex
\begin{document}
% Seleciona o idioma do documento 
\selectlanguage{brazil}
% ----------------------------------------------------------
% DADOS DO SEU TRABALHO
% ----------------------------------------------------------
\input{DadosDoTCC} % Preencha esse arquivo com seus dados!!!
% ----------------------------------------------------------
% PREÂMBULO
% ----------------------------------------------------------
\include{Textos/Preambulo/CapaTCC} % Capa do TCC
\include{Textos/Preambulo/FolhaAprovacao} % Folha de aprovação% 
% ----------------------------------------------------------
% ELEMENTOS PRÉ-TEXTUAIS
% Os nomes são auto-explicativos!
% ----------------------------------------------------------
\include{Textos/Agradecimentos} % Agredecimentos
\include{Textos/Resumo} % Resumo em português
\include{Textos/Abstract} % Resumo em inglês
\include{NaoAlterar/ListasBasicas} % Listas (Figuras, Quadros, e Tabelas) 
\include{Textos/ListaSiglas} % Lista de abreviaturas e siglas
\include{Textos/ListaSimbolos} % Lista de símbolos
% ----------------------------------------------------------
% SUMÁRIO - Não alterar!!!
% ----------------------------------------------------------
\pdfbookmark[0]{\contentsname}{toc}
\tableofcontents*
\cleardoublepage
% ---------------------------------------------------------
% ELEMENTOS TEXTUAIS
\textual \pagestyle{simple}
% 
% Agora seguem os capítulos principais!
% 
% Capítulo 1 - Introdução
\include{Capitulos/1_Introducao} % Introdução
% 
% Capítulo 2 - Objetivos
\include{Capitulos/2_Objetivos} % Objetivos (geral e específicos)
% 
% Capítulo 3 - Justificativas
\include{Capitulos/3_Justificativa} % Explica a relevância do trabalho
% 
% Capítulo 4 - Revisão Bibliográfica
\include{Capitulos/4_RevisaoBibliografica} % Fundamentos e estado da arte
% 
% Capítulo 5 - Metodologia e descrição das análises
\include{Capitulos/5_MateriaisMetodos}
%  
% Capítulo 6 - Resultados esperados e/ou parciais
\include{Capitulos/6_Resultados}
% 
% Capítulo 7 - Plano de trabalho e andamento
\include{Capitulos/7_Cronograma}

% ----------------------------------------------------------
% ELEMENTOS PÓS-TEXTUAIS
% ----------------------------------------------------------
\postextual \phantompart
% ----------------------------------------------------------
% Referências bibliográficas
\bibliography{Referencias}
% Apêndices - são textos elaborados pelo autor a fim de complementar sua argumentação. 
\include{Arquivos/Apendices}
% Anexos - são os documentos não elaborados pelo autor, que servem de fundamentação, comprovação ou ilustração, como mapas, leis, estatutos etc.
\include{Arquivos/Anexos}
% ----------------------------------------------------------
% FIM DO DOCUMENTO
% ----------------------------------------------------------
\end{document}

\end{document}